\section{The Learning Engine: Backpropagation}

The \emph{backpropagation algorithm}~\cite{Rumelhart1986} enables efficient gradient computation on the directed acyclic computational graph $\mathcal{G}$ representing the neural network.

\subsection{Forward and Backward Propagation}

The network transforms inputs $\vect{x} = \vect{a}^{(0)}$ through layers $l=1 \dots L$ via the forward pass:
\begin{align}
	\vect{z}^{(l)} &= \vect{W}^{(l)} \vect{a}^{(l-1)} + \vect{b}^{(l)}, \quad \vect{a}^{(l)} = \phi^{(l)}(\vect{z}^{(l)}) \\
	\loss &= \mathcal{L}(\vect{a}^{(L)}, \vect{y})
\end{align}
To update parameters $\params = \{\vect{W}^{(l)}, \vect{b}^{(l)}\}$, we compute gradients using the chain rule. Define the \emph{error signal} as $\vect{\delta}^{(l)} = \frac{\partial \loss}{\partial \vect{z}^{(l)}}$. The backward pass computes these recursively from $L$ down to $1$:

\begin{equation}
	\vect{\delta}^{(l)} = 
	\begin{cases} 
		\frac{\partial \loss}{\partial \vect{a}^{(L)}} \hadamard \phi'^{(L)}(\vect{z}^{(L)}) & \text{if } l=L \\
		\left[(\vect{W}^{(l+1)})^T \vect{\delta}^{(l+1)}\right] \hadamard \phi'^{(l)}(\vect{z}^{(l)}) & \text{if } 1 \le l < L
	\end{cases}
	\label{eq:backprop_recurrence}
\end{equation}

Once $\vect{\delta}^{(l)}$ is obtained, the gradients for layer $l$ are:
\begin{equation}
	\frac{\partial \loss}{\partial \vect{W}^{(l)}} = \vect{\delta}^{(l)} (\vect{a}^{(l-1)})^T, \quad
	\frac{\partial \loss}{\partial \vect{b}^{(l)}} = \vect{\delta}^{(l)}
\end{equation}

